% Irodov Part IV §4.2 94--109. Source user PDF pp.176--178;
% printed pp.167--169, source answers pp.317--318.
\electricoscproblem{94}{167}
铜金属薄板内自由电子整体沿板面法线微小平移 $x$ 后，表面产生电荷并形成静电回复场，引起电子等离子体振荡。铜中自由电子数密度 $n=0.85\times10^{29}\,\mathrm{m^{-3}}$，求其固有角频率。
\exsolution
离子晶格固定。电子整体位移 $x$ 在两表面形成电荷面密度大小 $\sigma=ne|x|$，内部净电场为 $E=\sigma/\varepsilon_0$，方向使电子回复。
单电子方程
$m_e\ddot x=-eE=-ne^2x/\varepsilon_0$，故
\[
\boxed{\omega_p=\sqrt{\frac{ne^2}{\varepsilon_0m_e}}
\simeq1.65\times10^{16}\,\mathrm{s^{-1}}}.
\]
\excheck 数密度单位必须为 $\mathrm{m^{-3}}$；该版底本纯文本把它识别为 $\mathrm{m^{-1}}$，与数密度定义不符。

\electricoscproblem{95}{168}
无电阻 LC 电路电容为 $C$、线圈自感为 $L$，持续自由无阻尼电振荡，电容两端电压最大幅度为 $V_m$。求任一时刻回路电流 $I$ 与电容电压 $V$ 的关系，分别采用电路方程和能量守恒推导。
\exsolution
回路方程 $L\dot I+V=0$，并有 $I=\pm C\dot V$（取一致的方向）。由消元得
$\ddot V+\omega_0^2V=0$，$\omega_0=1/\sqrt{LC}$。
若 $V=V_m\cos(\omega_0t+\alpha)$，
则 $I$ 与 $V$ 相差四分之一周期，
$I=\pm CV_m\omega_0\sin(\omega_0t+\alpha)$。
消去相位：
\[
\boxed{V^2+\frac LC I^2=V_m^2}.
\]
另由守恒
$\tfrac12CV^2+\tfrac12LI^2=\tfrac12CV_m^2$，也得到同一椭圆相轨迹。
\excheck $|I|\le V_m\sqrt{C/L}$；电容电压达到极值时电流恰为零。

\electricoscproblem{96}{168}
无电阻 LC 电路电容为 $C$、线圈自感为 $L$，初态开关断开，电容充至 $V_m$；$t=0$ 合闸。求：(a) $I(t)$；(b) 电容电场能与线圈磁场能相等瞬间的自感电动势。
\exsolution
设 $\omega_0=(LC)^{-1/2}$，
初态 $q(0)=CV_m$、$I(0)=0$。
电荷 $q=CV_m\cos\omega_0t$，
选放电方向为正，则
\[
\boxed{I(t)=V_m\sqrt{\frac CL}\sin\omega_0t}.
\]
电场与磁场能相等：
$CV^2/2=LI^2/2$，各占总能的一半，
故 $|V|=V_m/\sqrt2$。
依据无电阻 KVL，自感电动势与电容电压等模，
\[
\boxed{|\mathcal E_L|=\frac{V_m}{\sqrt2}}.
\]
\excheck 必须带 $1/\sqrt2$；该版书末公式文字抽取容易漏根号。

\electricoscproblem{97}{168}
无电阻 LC 振荡回路中，平行板电容器两极板被\emph{缓慢}拉开，使振荡固有频率提高至原来的 $\eta$ 倍；初始储能为 $W$。求移动极板所做功。
\exsolution
若极板分离远慢于振荡周期，系统为缓慢变参数谐振子。电容降低使
$\omega=(LC)^{-1/2}$ 增大。
对无耗散缓慢谐振，作用量
$J=W/\omega$ 为绝热不变量，所以
$W_f/W_i=\omega_f/\omega_i=\eta$，
由能量守恒，移动极板做功
\[
\boxed{A_{\rm adiabatic}=(\eta-1)W}.
\]
\textbf{书末条件差异：}该版书末给
$A=(\eta^2-1)W$。若额外要求在移动过程中电容电荷\emph{振幅固定}，则
$W\propto1/C\propto\omega^2$，确实出现该式；但无阻尼自由 LC 系统缓慢改变参数时电荷振幅并不固定。
\excheck 这两个结果属于\emph{不同的外加控制模型}。本稿保留原题“slowly”的通常绝热解释，并把书末差异列入勘误，而非无说明地替换。

\electricoscproblem{98}{168}
图4.27中电感 $L=2.5\,\mathrm{mH}$，两个电容 $C_1=2.0\,\mathrm{\mu F}$、
$C_2=3.0\,\mathrm{\mu F}$。两电容初始均充至同一电压 $V=180$V，合上中间电感支路开关后，求：(a) 固有振荡频率；(b) 电感电流最大值。
\exsolution
按原图两电容均跨接上下两公共节点，为并联，
$C_\Sigma=C_1+C_2=5.0\,\mathrm{\mu F}$。
总回路无电阻，固有角频率
$\omega_0=1/\sqrt{LC_\Sigma}$，
周期和普通频率为
\[
\boxed{T=2\pi\sqrt{LC_\Sigma}
\simeq0.70\,\mathrm{ms}},\qquad
\boxed{\nu_0=\frac1{2\pi\sqrt{LC_\Sigma}}
\simeq1.42\,\mathrm{kHz}}.
\]
初始电能全转为最大磁能：
$\tfrac12C_\Sigma V^2=\tfrac12LI_m^2$，
所以
\[
\boxed{I_m=V\sqrt{\frac{C_\Sigma}{L}}
\simeq8.0\,\mathrm A}.
\]
\excheck 书末(a)答的是周期 $T$，而题面只写“natural oscillation frequency”；本稿同时给出频率与周期，避免以毫秒冒充 Hz。

\electricoscproblem{99}{168}
图4.28的回路含两只相同电容 $C$、一只电感 $L$ 和开关，忽略电阻。左电容初始电压 $V_0$，右电容初始不带电；$t=0$ 合闸。求两电容各自电压随时间变化。
\exsolution
两电容经上方电感和下方导线耦合；以图中相同的上下极板参考方向记电压为 $V_1,V_2$。
两个电容上的总电荷守恒
$C(V_1+V_2)=CV_0$。
电感回路满足
$L\dot I=V_1-V_2$，并且
$I=-C\dot V_1=C\dot V_2$（选定方向可使两式同步）。
对差量 $u=V_1-V_2$，
$\ddot u+[2/(LC)]u=0$，
初值 $u(0)=V_0$、$\dot u(0)=0$。
故
\[
\boxed{V_1(t)=\frac{V_0}{2}(1+\cos\Omega t)},\quad
\boxed{V_2(t)=\frac{V_0}{2}(1-\cos\Omega t)},\quad
\boxed{\Omega=\sqrt{\frac2{LC}}}.
\]
\excheck $V_1+V_2=V_0$，初态分别为 $V_0,0$，半周期后两电容互换充电。

\electricoscproblem{100}{168}
无电阻 LC 电路电感为 $L$、电容为 $C$；线圈原处于恒定外磁场，穿过全部线圈的总外磁链为 $\Phi$。$t=0$ 将外磁场以远短于固有周期的时间撤去，求以后电路电流 $I(t)$。
\exsolution
外磁链骤降 $\Delta\Phi_{\rm ext}=-\Phi$，
根据短时间积分的法拉第—电感方程
$L\Delta I=-\Delta\Phi_{\rm ext}=\Phi$，
故撤场结束后初始电流
$I(0^+)=\Phi/L$。
撤场之后为自由 LC 振荡，
满足
$\ddot I+\omega_0^2I=0$，
并且电容原无电压，故
$\dot I(0^+)=0$。
于是
\[
\boxed{I(t)=\frac{\Phi}{L}
\cos\left(\frac{t}{\sqrt{LC}}\right)}.
\]
\excheck 磁链而不是单匝磁通入式；
撤场时间极短使电容电荷在该阶段几乎不变。

\electricoscproblem{101}{169}
电容电压按
$V(t)=V_me^{-\beta t}\cos\omega t$
作衰减振荡。求电压模在何时：(a) 与包络幅度相等；(b) 达到真正的局部极值。
\exsolution
(a) 电压模等于当前包络
$V_me^{-\beta t}$ 当且仅当
$|\cos\omega t|=1$，
所以
\[
\boxed{t_n^{\rm env}=\frac{n\pi}{\omega}},
\quad n=0,1,2,\ldots
\]
(b) 局部极值需 $\dot V=0$；
求导
$\dot V=-V_me^{-\beta t}
(\beta\cos\omega t+\omega\sin\omega t)$，
因此
\[
\boxed{\tan\omega t_n=-\frac\beta\omega},
\qquad
\boxed{t_n^{\rm ext}=
\frac{n\pi-\arctan(\beta/\omega)}{\omega}},
\]
选使 $t_n\ge0$ 的整数 $n$。
\excheck 因包络随时间衰减，电压实际局部极值一般不在包络接触点；初始 $t=0$ 是定义域端点。

\electricoscproblem{102}{169}
电阻为 $R$、电感为 $L$、电容为 $C$ 的串联 RLC 回路开关原断开，电容带初始电压 $V_0$，$t=0$ 合闸，开始欠阻尼振荡。求刚合闸时实际电压与位移振荡包络峰值之比。
\exsolution
令 $\beta=R/(2L)$、
$\omega_0=1/\sqrt{LC}$、
$\omega=\sqrt{\omega_0^2-\beta^2}$。
初始电流为零，即 $\dot q(0)=0$。
因此
$q(t)=CV_0e^{-\beta t}
[\cos\omega t+(\beta/\omega)\sin\omega t]$。
括号的相位合成振幅为
$\sqrt{1+(\beta/\omega)^2}=\omega_0/\omega$，
电容电压的指数包络峰值为
$V_m=V_0\omega_0/\omega$。
故
\[
\boxed{\frac{V_0}{V_m}
=\frac{\omega}{\omega_0}
=\sqrt{1-\frac{R^2C}{4L}}}.
\]
\excheck 欠阻尼条件为 $R^2C<4L$，根号非负；不得漏掉平方根。

\electricoscproblem{103}{169}
串联 RLC 回路电感 $L$、电容 $C$，自发自由阻尼电流
$I(t)=I_me^{-\beta t}\sin\omega t$。求电容电压 $V_C(t)$，特别是 $t=0$ 的电压。
\exsolution
回路方程
$L\dot I+RI+V_C=0$。
自由阻尼参数
$R=2L\beta$、
$\omega_0^2=\omega^2+\beta^2=1/(LC)$。
代入 $I(t)$，选电流与电容端电压的对应有向关系：
\[
\boxed{V_C(t)=
-LI_me^{-\beta t}
[\omega\cos\omega t+\beta\sin\omega t]}.
\]
写相位式，其振幅
$L I_m\sqrt{\omega^2+\beta^2}
=I_m\sqrt{L/C}$；
初始电压
\[
\boxed{|V_C(0)|=LI_m\omega
=I_m\sqrt{\frac{L/C}{1+\beta^2/\omega^2}}}.
\]
\excheck 不同电流箭头定义会使 $V_C$ 整体反号，但电压模与相位差保持不变。

\electricoscproblem{104}{169}
串联 RLC 电路 $C=4.0\,\mathrm{\mu F}$、
$L=2.0\,\mathrm{mH}$、$R=10\,\Omega$。
在自由阻尼振荡的某个电流极大时刻，求磁场能与电容电场能之比。
\exsolution
电流取极大瞬间 $\dot I=0$，
回路 KVL
$L\dot I+RI+V_C=0$
简化为
$V_C=-RI$。
故
\[
\frac{W_L}{W_C}
=\frac{LI^2/2}{CV_C^2/2}
=\frac{L}{CR^2}.
\]
代入得
\[
\boxed{\frac{W_L}{W_C}
=\frac{2.0\times10^{-3}}
{(4.0\times10^{-6})(10)^2}=5.0}.
\]
\excheck 电流最大时 $\dot I=0$，但\emph{电容电压并非零}；阻尼回路与无损 LC 的能量关系不同。

\electricoscproblem{105}{169}
串联 LC 振荡电路中，两只线圈的电感分别为 $L_1,L_2$、电阻分别为 $R_1,R_2$，忽略相互电感。希望用一只线圈替代它们，使振动频率与品质因数不变。求新线圈参数。
\exsolution
串联相同电流通过两线圈，磁链相加，
$L_{\rm eq}=L_1+L_2$，电阻亦串联相加
$R_{\rm eq}=R_1+R_2$。
电容 $C$ 保持不变时，
自由振动角频率
$\omega_d^2=1/(LC)-(R/2L)^2$，
品质因数由 $L,R,C$ 决定；
采用等效参数自然使两个量相同。
\[
\boxed{L=L_1+L_2,\qquad R=R_1+R_2}.
\]
\excheck 假设两线圈互感可忽略，
否则串联同向或反向耦合还需 $\pm2M$。

\electricoscproblem{106}{169}
LC 阻尼电路品质因数 $Q=5000$，振荡频率
$\nu=2.2\,\mathrm{MHz}$。求电流\emph{振幅}下降 $\eta=2.0$ 倍所需时间。
\exsolution
弱阻尼自由电流幅度
$I_m(t)=I_m(0)e^{-\beta t}$，
且 $Q\simeq\omega/(2\beta)$、
$\omega=2\pi\nu$，
所以
$\beta\simeq\pi\nu/Q$。
于是
\[
\boxed{t=\frac{\ln\eta}{\beta}
\simeq\frac{Q\ln\eta}{\pi\nu}}.
\]
代数值
\[
\boxed{t\simeq\frac{5000\ln2}{\pi(2.2\times10^6)}
=5.0\times10^{-4}\,\mathrm s=0.50\,\mathrm{ms}}.
\]
\excheck 该版书末的公式符合此计算，但数值提取似为“0.5 s”而非毫秒；代题面 MHz 时必为约0.5 ms。

\electricoscproblem{107}{169}
RLC 回路 $C=10\,\mathrm{\mu F}$、
$L=25\,\mathrm{mH}$、$R=1.0\,\Omega$。
求电流振幅降低 $e$ 倍需经历多少个\emph{实际阻尼振荡周期}。
\exsolution
$\beta=R/(2L)$、
$\omega_d=\sqrt{1/(LC)-\beta^2}$。
幅度降低 $e$ 倍需 $t=1/\beta=2L/R$，
期间周期数
\[
N=\frac t{2\pi/\omega_d}
=\frac{\omega_d}{2\pi\beta}
=\frac1{2\pi}\sqrt{\frac{4L}{CR^2}-1}.
\]
代入
\[
\boxed{N\simeq
\frac{\sqrt{4(0.025)/(10^{-5})-1}}{2\pi}
\simeq16}.
\]
\excheck 此处 $N$ 是周期数，不是电荷的极值次数，不能额外乘2。

\electricoscproblem{108}{169}
RLC 自由阻尼振荡电路品质因数为 $Q=5.0$，
求实际振荡角频率 $\omega$ 与无阻尼固有角频率 $\omega_0$ 的相对差百分比。
\exsolution
若品质因数按阻尼周期的 $\omega/(2\beta)$ 定义，
$Q=\omega/(2\beta)$，
则
$\omega_0^2=\omega^2+\beta^2
=\omega^2(1+1/(4Q^2))$，
因此
\[
\boxed{\frac{\omega_0-\omega}{\omega_0}
=1-\frac1{\sqrt{1+1/(4Q^2)}}}.
\]
$Q=5$ 时约
$\boxed{0.50\%}$。
\excheck 使用 $Q=\omega_0/(2\beta)$ 的另一常用定义也可在高 Q 下得到约0.5%，区别在更高阶修正。

\electricoscproblem{109}{169}
图4.29：电源电动势 $\mathcal E=2.0$V，
内阻 $r=9.0\,\Omega$，电容
$C=10\,\mathrm{\mu F}$，
电感 $L=100\,\mathrm{mH}$，
线圈电阻 $R=1.0\,\Omega$。
原先开关闭合达稳态，某刻断开，电容—线圈回路自由振荡。求断开后：(a) 初始振荡能量；(b) $t=0.30$s 的剩余能量。
\exsolution
闭合稳态电感短路，初始电流
$I_0=\mathcal E/(r+R)$，
电容电压 $V_{C0}=RI_0$。
因电容与线圈在稳态时为并联，
初始振荡能量
\[
W_0=\frac12LI_0^2+\frac12CV_{C0}^2
=\frac{\mathcal E^2(L+CR^2)}{2(r+R)^2}.
\]
数值
$\boxed{W_0\simeq2.0\,\mathrm{mJ}}$。
弱阻尼下能量包络衰减
$W(t)\simeq W_0e^{-Rt/L}$，
所以
\[
\boxed{W(0.30\,\mathrm s)\simeq
(2.0\,\mathrm{mJ})e^{-3}
\simeq0.10\,\mathrm{mJ}}.
\]
\excheck $r$ 在开关断开后被排除在振荡回路外，后续衰减取线圈电阻 $R$；能量指数衰减常数为 $R/L$ 而不是电流幅度的 $R/(2L)$。
